\scanpage{538}%
\mspage{449}%
% archive sheet number: 537
de $\alpha$ équivaut à celle d'un élément de $\mathop{\unsure{\underline{\lim}|\varinjlim}}\limits_{i\in I}\operatorname{Hom}(z,x_i)$, \deleted{\unsure{cet}} est donc donné par un

\[\alpha_{i_0}:z\longrightarrow x_{i_0},\]

la donnée de $\beta$ équivaut à celle d'un élément de $\varprojlim_{i\in I}\operatorname{Hom}(x_i,z)$, i.e. de flèches

\[\beta_i:x_i\longrightarrow z,\]

avec la condition de compatibilité habituelle. \inserted{Quitte à remplacer $I$ par $I'$, avec $I'\to I$ cofinal,} OPS $I$ ordonné\inserted{(et grd devant $\pi_0$, si $I$ initial l'était)}, et la condition \deleted{\unsure{de}} de compatibilité sur les $\beta_i$ s'écrit

\[\beta_i\tau_{ji}=\beta_j\qquad\text{si }i\leq j\text{ ds }I,\]

où $\tau_{ji}:x_i\mathrel{\unsure{\relbar|\to}}x_j$ désigne la flèche de transition. \deleted{\unsure{D|Il}}Il faut expliciter, en termes de \unsure{nos|ces|ses} données, les conditions

\[\alpha\beta=\mathrm{id}_z,\qquad\beta\alpha=\mathrm{id}_{\deleted{\unsure{z|x}}\underline{x}}\unsure{,|!}\]

\diagram{\begin{tikzcd}[ampersand replacement=\&] z\arrow[r,"{\alpha_{i_0}}"]\arrow[dr] \& x_{i_0}\arrow[r,dash,"{\tau_{i,i_0}}"]\arrow[d,dash,"{\beta_{i_0}}"] \& x_i\arrow[dl,dash,"{\beta_i}"']\\ \& z \& \end{tikzcd}}

La relation $\alpha\beta=\mathrm{id}_z$ équivaut à \deleted{\unsure{$\beta$}}

\begingroup\makeatletter\tagsleft@true\makeatother
\[\deleted{\unsure{\beta z|\alpha\beta}}\ \beta_{i_0}\alpha_{i_0}=\mathrm{id}_z,\tag*{$(*_{i_0})$}\]
\endgroup

elle implique, d'ailleurs, en posant pour $i\geq i_0$

\[\alpha_i\overset{\mathrm{déf}}{=}\tau_{i,i_0}\alpha_{i_0}:z\longrightarrow x_i,\]

que l'on a pour $i\geq i_0$

\begingroup\makeatletter\tagsleft@true\makeatother
\[\beta_i\alpha_i=\mathrm{id}_z\qquad\forall i\geq i_0.\tag*{$(*_i)$}\]
\endgroup
